\chapter*{Lời Mở Đầu}
\addcontentsline{toc}{chapter}{Lời Mở Đầu}
\markboth{Lời Mở Đầu}{Lời Mở Đầu}

\begin{truyen}{Lời Mở Đầu}{Opening Words}
\chuhoa{C}{uộc đời có những quy luật riêng — ai hiểu sớm thì đỡ phải trả giá.}
\textit{(Life has its own laws --- those who understand them early pay fewer prices.)}

Quyển sáu mang đến những câu chuyện về vòng quay của cuộc đời: nhân quả, bản sắc, buông bỏ và những người thầy vô hình xuất hiện đúng lúc ta cần.
\textit{(Volume six brings stories about the cycle of life: karma, identity, letting go, and the invisible teachers who appear exactly when we need them.)}

Đôi khi bài học quý nhất không đến từ lớp học, mà từ một người gặp ngẫu nhiên, từ một thất bại đúng lúc, từ một câu nói vô tình mà ta nhớ mãi.
\textit{(Sometimes the most valuable lessons do not come from the classroom, but from a chance encounter, a timely failure, an accidental sentence we remember forever.)}

Hãy đọc quyển này với tâm thế rộng mở — vì bài học có thể đến từ bất kỳ câu chuyện nào, bất kỳ trang nào.
\textit{(Read this volume with an open mind --- because lessons can come from any story, any page.)}
\end{truyen}

\begin{baihoc}
Cuộc đời không vận hành theo ý muốn của ta --- nhưng ta có thể học cách vận hành cùng với nó.
\textit{(Life does not operate according to our wishes --- but we can learn to operate alongside it.)}
\end{baihoc}

\begin{tcolorbox}[colback=truyenblue!6,colframe=truyenblue,coltitle=white,fonttitle=\bfseries,title={Easy English},arc=4pt,boxrule=1pt,breakable]
\textbf{Short English to remember:}

Life has a rhythm --- learn to move with it.

Invisible teachers are everywhere.
\end{tcolorbox}

\begin{tcolorbox}[colback=truyengold!10,colframe=truyengold!70!black,coltitle=black,fonttitle=\bfseries,title={Tự Học Nhanh},arc=4pt,boxrule=1pt,breakable]
1. Trong cuộc đời em cho đến nay, điều gì em coi là nhân quả rõ ràng nhất? \textit{In your life so far, what do you see as the clearest example of cause and effect?}

2. Điều gì em đang cần buông bỏ nhưng chưa làm được? \textit{What do you need to let go of but have not been able to?}

3. Ai là ``người thầy vô hình'' có ảnh hưởng lớn nhất đến em? \textit{Who is the "invisible teacher" with the greatest influence on you?}
\end{tcolorbox}

\ngancach
